%!TEX program = lualatex
% SPDX-License-Identifier: CC-BY-SA-4.0
% latex-atlas 02: marginal scholia. Isolates sidenotes, catchword, folio.
% Geometry gives a broad outer margin; marginpar holds the gloss.
\documentclass[11pt,twoside]{article}

\usepackage{iftex}
\ifLuaTeX\else
  \PackageError{latex-atlas}{This example requires LuaLaTeX}
    {Run lualatex, not pdflatex.}
\fi

\usepackage{fontspec}
\usepackage{xcolor}
\usepackage{geometry}
\usepackage[tracking=true,protrusion=false,expansion=false]{microtype}

\defaultfontfeatures{Renderer=HarfBuzz}
\IfFontExistsTF{EB Garamond}{
  \setmainfont{EB Garamond}[
    Numbers={OldStyle,Proportional},
    Ligatures={TeX,Common,Rare}
  ]
}{
  \setmainfont{TeX Gyre Pagella}[
    Numbers={OldStyle,Proportional},
    Ligatures={TeX,Common}
  ]
}

% Paper tint is per-example tuning for this wide-margin page, not a shared
% token yet; the three examples deliberately differ slightly (see README).
\definecolor{parchment}{HTML}{F4EBDD}
\definecolor{ink}{HTML}{231F1A}
\definecolor{glossink}{HTML}{5A4A3A}
\pagecolor{parchment}
\color{ink}
\geometry{
  a4paper,
  top=28mm, bottom=32mm,
  inner=25mm, outer=52mm,
  marginparwidth=34mm, marginparsep=6mm
}
\setlength{\parindent}{1.1em}
\setlength{\parskip}{0pt}
\pagestyle{plain}

% A gloss in the margin: small, ragged, in a quieter ink.
\newcommand{\gloss}[1]{\marginpar{\footnotesize\raggedright\color{glossink}#1}}

% Catchword: first word of the next page, set alone at the foot.
\newcommand{\catchword}[1]{%
  \vfill
  \begin{flushright}
    {\small\textit{#1}}\par
    {\scriptsize\textls[80]{\textsc{catchword}}}
  \end{flushright}
}

\begin{document}

\begin{center}
  {\Large\textls[70]{\textsc{Of Pendulums}}}\par
  \vspace{0.3em}
  {\small\textit{with the gloss of a later reader}}
\end{center}

\vspace{1em}

A heavy ball hung by a thread, and drawn aside from the vertical,
returns and passes beyond it.\gloss{The isochrony holds only for small
arcs; at wide excursions the period lengthens.} For small excursions the
time of one swing depends upon the length alone, and not upon the weight
nor the width of the arc.

The length is reckoned from the point of suspension to the centre of the
ball.\gloss{Measure to the centre, not to the top of the ball; the thread
itself is taken as without weight.} A change of a finger's breadth in a
line of three feet sensibly alters the going of a clock.

Whence the pendulum serves to divide the day into equal parts, and to
try the figure of the earth.\gloss{Huygens, \textit{Horologium}, 1673;
Richter's Cayenne retardation, 1672.} What is here set down of one ball
holds of every ball, provided the thread be fine and the arc narrow.

\catchword{Of falling}

\vfill
\begin{center}
  {\footnotesize\textit{--- 7 ---}}
\end{center}

\end{document}
